前两章分别从显式几何法则约束以及局部底层离散形函数畸变这两个切面层级，剥析了物理模拟体系中由局部高频硬核干扰所引爆的时间积分步缩短与空间离散矩阵条件数病态难题，并成功演示了如何通过改写、定制系统表示与解构空间底层映射，自源头将不良数值刚度彻底剥离以逆转乾坤的解题哲思。但是，当问题的维度一再极速攀升并扩张至包裹着数千万乃至于数亿量级超级自由度节点的结构动力学与超巨型系统模态特征分析疆域时，阻滞性能进一步狂飙的困难重心，已悄然间发生了极具系统性维度的抽象转移。此时此刻：哪怕系统网格里连一个畸形的恶质单元都找不出来，单单是要从数以千万计的巨量方程交织海中，以极度高昂的效率提取并拼接出一套足以严丝合缝地囊括、张罗并反映全系总成中至关重要之"软"态系统低频整体大位移、大振荡波长面貌属性的极致完美缩减宏观响应子空间，就已然演化成为了限制计算机极致算力的那道不可跨越的浩瀚天堑与最致命瓶颈。这便是本章顺延前序理论所挺进的全篇最后一层，也是最为庞观宏远、聚焦于巨型代数算盘谱逼近效率不平衡症结——广义大规模系统特征求值中的频域频段表达严重病态与降阶维度失衡问题。

\section{引言}

通过捕捉系统的关键低阶频段特征谱及其形变矩阵，以求解并刻画出掌控物理运动根本轨迹行为的特征分析方程体系，长久以来都是现代超尖端精密工程开发验证乃至于逼真物理图形仿真动画控制的核心枢纽基石。在当前海量多层装配模型的逼临情境里，全盘包揽式的宏观统筹解算（如调用基于高性能 Krylov 矩阵降秩理念的框架 SLEPc 混搭硬核直接多节点分解工具 MUMPS 之流），虽然能产出精度绝不打折扣的标准金律，却因为躲不开病态条件数积累引发的极端迟钝迭代、和令人难以负担的分布式节点海量矩阵重组计算代价导致庞大的"存储阻塞墙"效应而遭遇严重可伸缩性能卡壳。更为致命的是，这种全系统统包的求解算器骨子里即丧失了在极度严苛工程环境下用于多频次修改更新的设计迭代灵敏弹性：即便是模型边缘一个螺丝孔几何的微改动，也会诱发不可逆的全网格分解连锁惩罚运算，令整个分析流转效率低得不堪重用。

鉴于这种窘境，主打子空间拼接哲理且支持化整为零区域解体的"模态综合法"（Component Mode Synthesis，CMS）在航空造机等庞然大物分析流程及商用老牌分析重器中依然稳居统治之位。其大核价值正在于通过区域分割法巧妙促成了针对局部组件更新重新解析与知识产权屏蔽的可重用大业。然而，经典 CMS 以极大的收敛性与数值精确性折损，换取了计算敏捷性空间，这个残酷的妥协，恰恰成为了触发更深维度病态谱失衡难题的渊薮：\textcolor{IsConclusion}{其\textbf{收敛速度往往不尽如人意}。}

\subsection{相关工作}

大规模工程结构低频模态的模态分析是一个典型的特征值问题。标准的"整体式"方法将特征求解器直接应用于完整的、未分区的系统矩阵。典型的求解器包括 Krylov 子空间方法（如 Lanczos 或 Arnoldi 方法）\cite{Hernandez2005SLEPc, Qiu2015Spectra}、Jacobi-Davidson (JD) 方法 \cite{Sleijpen2000JD} 以及局部最优块预处理共轭梯度法 (LOBPCG) \cite{knyazev2001lobpcg}。随着模型规模的增加，这些方法的计算成本变得难以承受。然而，它们不仅产生巨大的内存开销，更重要的是在迭代过程中需要重复进行线性求解（例如用于位移-逆变换或预处理步骤）。

\subsubsection{模态综合法}

这一计算瓶颈促使了基于区域分解的技术的发展，例如模态综合法（Component Mode Synthesis, CMS） \cite{seshu1997substructuring, William1985CMS, bourquin1992component}。具有里程碑意义的 Craig-Bampton (CB) 方法 \cite{craig1968coupling} 几乎被广泛应用于所有商业有限元分析软件 \cite{Vysok2017CurrentCO}。它从两组模态构造缩减基：固定界面正规模态和静态约束模态。然而，对于现代高保真有限元模型，CB 方法面临严重的扩展性瓶颈。随着子结构特征模态数量的增加，其计算变得非常缓慢，尤其是当网格规模较大时。此外，界面的自由度数量将变得非常庞大 \cite{KRATTIGER2019Riview}。

\subsubsection{界面模态缩减法}

鉴于界面系统规模庞大且稠密，通常通过特征分解对其进行全局或局部进一步缩减。
\citet{Castanier20101Characteristic} 提出在界面全局组装后，直接求解界面凝聚系统矩阵的特征值问题。这种方法被称为系统级缩减，效率较低。当界面简单且可分离为若干独立界面时，全局 Schur 补可以在这些独立界面及其相邻子结构上局部计算 \citep{Cheon2023Enhanced,KRATTIGER2019Riview}。当界面复杂且涉及多个交叉点时，特征分解应用于子结构级别的凝聚界面矩阵 \cite{HONG2013403, Kuether2017Model}。此外，\citet{BATTIATO2018187} 提出了 Gram-Schmidt 界面模态方法来处理非协调网格问题。

\subsubsection{划分与区域分解法}

与区域分解法 (DDM)~\cite{Smith1997DDM} 类似，模态综合法采用分治策略将计算域划分为多个不重叠的子域。存在多种划分策略：基于网格空间坐标的几何划分、利用主成分分析 (PCA) 对顶点进行聚类~\cite{HUANG2006927}、或依赖图划分算法如 METIS~\cite{Karypis1995graph, Karypis1997METIS} 和 CHACO~\cite{Hendrickson2011}。在本工作中，我们不专注于优化划分策略，而是采用基于网格几何的简单几何划分方法。

\subsection{本章方法与贡献}

回溯经典传统 CMS （尤其是工业标杆级的 Craig-Bampton 方程论）实施体系，这头系统巨兽被大刀阔斧分割成了多块小区域子结构后，每一区块便开始在固定僵死的强力边界约束束缚下，孤立无援地剥离运算着只属于其自身地盘的一套截断基模组合。表观面上，这一切看似制造出了一系列相互解耦互不制肘可全火力并行拉动降维推演的最优局部降维基底；然而深入基向量拼叠合成频域解析空间时却骇然发现，这道硬性加诸的死板固定边界指令就像一道铁锁，残酷而直接地掐断了局部提取基函数跨越拼接交界点之际用以吻合连接并还原系统整体长波低频震荡特征的核心命脉——相位连续性。

由于这股极其病态的"相位残缺缺陷"，传统解法只能像填鸭般疯狂补充截留并压榨进大量原本并不该出现的极高波数局域高频额外基函数；更为悲剧的是，这股无奈吞入的高频噪音不仅和真实欲求的大宏观缓和低频空间模态近似正交排斥、在精度补充上几乎属于隔靴搔痒徒劳无功，反而更进一步造成降级表达中的低效率。与之结伴而行的是，为了使得相邻的组件不致分裂并强行粘合这股撕裂断层，算法不得不痛苦而又死板地依赖一套计算资源消耗极度可怕的所谓界面边界约束模态生成法：特别是在所需抽取的特征波模数量尚属微少精简的大型并行环境中，组装乃至储存与之交织衍生的界面 Schur 补等超稠密高代价算子体系不仅成了制约内存爆满的二次方量级灾难梦魇，更俨然变成阻塞通讯并疯狂拖垮降级扩展速率的最核心瓶颈。

为了从这个死循环病灶困局中解脱并矫正其特征表征缩减域失衡的顽疾，本章将目标直接锚定于放弃于旧有固化划分的子体系池底中去修修补补添加高频累赘的挣扎，转向以颠覆与重定义底层降阶拼图框架构造空间为其根本法门，旨在催生出一股更能自然平滑捕捞收容系统大尺度缓波柔性能量流的全频响应新组合体系。

我们本质上带来了一套针对降阶特征空间体系重塑与重造内涵逻辑：绝不囿限于某组僵死的局部划分块状集，而是创意性大开地引进了基于\textbf{多重空间层叠并错位交织布置的划分}法则。借助于多套相互错落切割系统所滋生出的自然频域位置更迭特性，不同切割子域内的局部模构响应在其交相重叠处能以完美的相位互补补齐彼此原本各自死死锁缺的部分相位图景谱线，如此一来便极其顺畅地彻底规避并减少了对于额外补集高频扰动段态空间的巨量无理索取，直接建构拼图出了一副更臻至完美圆融且具备高度捕获低能波表现特性的充实特征大系统模型空间。

同时，为斩断随之并生的昂贵界面耦合运算宿命毒瘤，本章充分巧用"这一组划分下的硬邦僵死交界带，在错开交织布置的另外那套划分组中极大概率只是一段处于柔性域内极其平顺安全的内域点"这一特异位置置换红利，提出了一种惊艳的完全不需要计算稠密 Schur 补亦不用再次经历高昂解方程提特征值的 IMR 取代公式逻辑，全面用其余区域推导传递转换过的系统边界静态局部响应反补估算算法，将过往天代价级的局限边界基函数处理操作，极其优雅地一并踢出了系统解算的整体链路中。

整体而言，本工作的研究贡献结晶囊括如下：
\begin{itemize}
\item 一套基于\textbf{错位多重网格划分融合策略}构架的大规模降维重组基准机制，经由自然调和互补缺失的截断面相位信息，令系统在获取提取宏大低阶逼近精度方面获得跃升级的精度赋能。
\item 创设发展出一组彻底颠覆陈规、完全跳脱组装并求解决庞大计算体量的\textbf{无 Schur 补及解绑局部特征演算的超高效 IMR 公式计算结构}，直接以静力局部重平衡方程取而代之解决了原本棘手的系统边际拼接昂贵瓶颈。
\item 执行并实施了详实透彻的大量数理检验环节，以对比之前传统 CMS 系列手法全面展示我们成果狂飙突破 3 个甚至更大量的精度进阶数量级提升跨度优势。本篇同时重点显明了这套理念重塑下带来的无解工业实战应用红利优势：（1）完美兼容极简修改模型网络特征节点触发高频小颗粒极小资源反向重算推演的快速修改循环；（2）在万级规模多超级超算核心集群与分布式复杂集群生态环境里极尽出色的通信低耗以及无上限庞大尺度计算扩展潜能。
\end{itemize}
